\section{Erasure}
\label{sec:erasure}

Everywhere else in the store a withdrawn claim stays readable beside its
retraction. A data subject who exercises a right to erasure, such as the one in
Article~17 of the GDPR~\cite{gdpr}, asks for the opposite. \op{erase} is the one
operation that deletes rows, and it is kept narrow enough that it cannot serve as
a general delete.

\subsection{What the operation does}

\shipped{}
\begin{itemize}[leftmargin=*,itemsep=2pt]
\item \textbf{Scope.} It removes the rows that name one entity, as the subject or
  as the target of an edge, in the caller's own (organization, project) file. A
  property whose scalar value happens to spell the entity's key is not a
  reference to it and is left.
\item \textbf{Authority.} It requires an administrator of the organization it
  acts in. A platform administrator acting inside another organization is
  refused unless that organization has made them one of its own administrators.
\item \textbf{Hold.} If any row that names the entity carries the litigation-hold
  bit, it removes nothing and reports how many rows are held. An erasure that
  skipped the held rows would report the subject erased while rows under hold
  still named it.
\item \textbf{Record first.} It is refused on a deployment with no durable audit
  store. The audit record is appended inside the deleting transaction, after the
  deletes and before the commit, so a failure to record rolls the deletion back.
  If the commit fails after the record was appended, a second record with an
  error outcome is appended under the same digest, since the trail is append-only.
\item \textbf{Bytes, not only rows.} The transaction runs with SQLite's
  \texttt{secure\_delete}, which overwrites freed pages; the full-text index is
  switched to its secure-delete mode, which removes the terms rather than masking
  them; and after the commit the write-ahead log is checkpointed and truncated,
  so frames written before the erasure leave the file~\cite{sqlitedocs}. A test
  reads the bytes of the database file, its write-ahead log and its rollback
  journal before and after, in both journal modes, and requires the subject's
  key to be present before and absent after. If the checkpoint fails, the
  erasure stands, every
  read has lost the rows, and the failure is logged: the bytes leave at the next
  checkpoint.
\item \textbf{Bound.} One call removes at most 10{,}000 rows; the receipt says
  how many still name the entity, and each further call has its own receipt.
\item \textbf{Reason.} The caller supplies a reason, which the audit trail keeps
  after the subject is gone, so a reason that contains the entity's key is
  refused.
\end{itemize}

\subsection{The receipt}

The receipt lists the content addresses of the removed rows in ascending order,
a digest over them computed as a content address is ($H$ over the
length-prefixed identifiers), the instant, the erasing identity, the sources the
removed rows were filed from, and a status for each store that may hold what the
organization says about the subject. The audit record carries the digest, the
count, the reason and the graph's namespace, and neither the entity nor any
removed field. A receipt is verified by recomputing the digest from its
identifiers and finding the audit record whose resource is that digest.

\paragraph{What a receipt establishes.} That an erasure removing exactly the
listed rows was recorded, at the stated instant, by the stated identity, in the
stated graph; and, with the test above, that the rows' bytes were no longer in
the database file once the checkpoint succeeded.

\paragraph{What it does not establish.}
\begin{itemize}[leftmargin=*,itemsep=2pt]
\item \emph{That nothing else holds the subject.} The receipt names two other
  stores in the same binary --- the knowledge pages, memories and sources with
  their vectors, and the \texttt{/v1/ai/memory} service --- whose rows carry no
  subject key and which offer no delete by subject. Both are listed as
  unsupported, so at \texttt{1fffdf408} every receipt reports
  \texttt{complete: false}, and those stores must be handled out of band. Backups,
  replicas, exports and anything a caller copied out of earlier reads are outside
  the operation entirely.
\item \emph{That the subject will stay erased.} The receipt lists the sources the
  rows were filed from. The store does not own those documents, and ingesting one
  again files the subject again.
\item \emph{That the identifiers reveal nothing.} A content address is a hash of
  low-entropy fields. Whoever holds a receipt and can guess an erased assertion
  in full can confirm the guess, so the receipt is itself personal data and is
  held as such. The audit trail keeps only the digest over all identifiers,
  which confirms a guess only of the entire erased set.
\end{itemize}

\subsection{Litigation hold}

Article~17(3)(e) of the GDPR exempts processing that is necessary for the
establishment, exercise or defence of legal claims~\cite{gdpr}, and a record
that is under hold may not be destroyed. The store models the hold as a bit on
each row, outside the content address because it describes the record rather
than the world. Erasure is all or nothing against it, and the deleting statement
itself also requires the bit to be clear.

At \texttt{1fffdf408}, no operation sets the bit. The refusal path is exercised by
a test that sets it directly in the database, and HIP-1198 records the same
gap~\cite{hip1198}. An operation that places and releases a hold, with its own
audit record, is required before the refusal can occur in service
\proposed{}. So is a disposal operation that expresses an organization's own
retention decision and that a hold outranks, which HIP-1198 names as the
condition for leaving its alpha stage \proposed{}.
